\section{具身智能的学术边缘史}

上一章说明语言和视觉的边界，本章进入学术史。王鹤强调，“具身智能”和“机器人”最初并不是同一拨学者的同一套叙事。具身智能更多起源于计算机视觉研究者对下一阶段问题的寻找：从 passive perception 走向能行动、能改变观察、能闭环学习的智能体。机器人传统学者则长期关注控制、机械结构、优化、工业场景和具体硬件。

这个区别解释了为什么具身智能早期最重要的任务不是灵巧操作，而是导航。导航让视觉研究者保持在相对熟悉的舒适区：agent 移动自身位置，摄像头看到新画面，环境本身不必被改动。它已经具备感知-动作循环，却比抓取、开门、接触和力控更容易标准化。因此 Object Goal Navigation、Point Goal Navigation、Habitat 等任务和平台成为早期叙事中心。

\begin{figure}[H]
\centering
\includegraphics[width=0.96\textwidth]{embodied-ai-origin.png}
\caption{具身智能学术边缘史：从视觉、导航、感知和动作循环走向主流。自制概念图，依据 00:25:08--00:41:15 对谈内容整理。}
\end{figure}

\begin{knowledgebox}{读图：学术流派不是从“机器人产业”自然长出来的}
图中从 Vision 到 Navigation，再到 Perception 和 Action，说明具身智能最早是视觉研究者对“下一代视觉任务”的扩展。它后来吸纳 robot learning、抓取、强化学习和大模型，但核心叙事最早是：视觉不能只被动识别，而要进入行动闭环。
\end{knowledgebox}

\subsection{Perception-Action Loop}

本节把具身智能的底层范式讲清楚。Perception-Action Loop 是“先通过感知决定行动，行动改变环境或自身位置，新的感知再指导下一步”。导航是最简单例子：机器人往前走，坐标和相机视角改变，于是看到新的环境。操作任务更复杂：抓取瓶子后，瓶子位置、手中状态和后续可行动作都改变。

与单步抓取相比，闭环更能体现具身智能的研究价值。早期许多抓取方法把问题变成点云或几何推理：看一眼物体，预测抓取点，执行一次，任务结束。它虽然发生在机器人上，但仍像一个三维模式识别任务。具身智能要追求的是多步闭环：每次行动都改变世界状态，模型需要重新观察、重新计划、继续执行。

\begin{figure}[H]
\centering
\includegraphics[width=0.96\textwidth]{perception-action-loop.png}
\caption{Perception-Action Loop：具身智能的核心叙事是感知、行动和反馈闭环。自制概念图，依据 00:25:08--00:41:15 对谈内容整理。}
\end{figure}

\begin{knowledgebox}{读图：闭环让视觉从分类变成行动}
从 Observe 到 Infer，再到 Act、Feedback 和 Adapt，关键变化是环境不再只是输入图片，而是会因为动作产生新状态。读图时要注意最后的 Adapt：具身智能不是一次预测，而是连续地用反馈修正策略。
\end{knowledgebox}

\begin{importantbox}{标志性事件}
王鹤提到，李飞飞将具身智能称为计算机视觉未来的三颗北极星之一，黄仁勋在 NVIDIA 场合强调下一代 AI 是 embodied AI。这类表述把学术边缘方向推向主流，也让视觉、机器人、强化学习和大模型社区开始围绕同一个词汇聚。
\end{importantbox}

\subsection{本章小结}

具身智能的学术史是一场“视觉主动化”的迁移：从图像识别转向导航，从单步推理转向闭环行动，再逐步吸纳 manipulation、robot learning 和大模型通用任务能力。

